\section{Data Introduction}

This project is built on \textbf{EdNet-KT4}, the largest publicly available
educational interaction dataset, released by Riiid AI Research and collected
from \emph{Santa}, a Korean TOEIC preparation platform. KT4 is the most
detailed of the four EdNet sub-datasets: it logs every fine-grained student
action — not only question responses, but also page navigation, audio/video
playback events, payments, and coupon redemptions. This level of detail makes
it a strong fit for a CS313 data-mining project: the full pipeline (cleaning,
integration, transformation, reduction, feature engineering, and modeling)
can be exercised on a real, large-scale, multi-modal log.

\subsection{Scale and provenance}

\begin{table}[H]
\centering
\small
\begin{tabular}{ll}
\toprule
Property & Value \\
\midrule
Total interactions   & 131{,}441{,}538 rows \\
Total users          & 297{,}915 unique students \\
Collection period    & 2018-08-27 to 2019-12-01 (461 days) \\
Source platform      & Santa (Riiid) — TOEIC AI tutor \\
Raw size             & 6.4 GB across 297{,}915 per-user CSV files \\
Consolidated parquet & 1.3 GB (snappy-compressed, single file) \\
\bottomrule
\end{tabular}
\caption{EdNet-KT4 at a glance.}
\end{table}

\subsection{Schema}

The interaction log has 8 columns. Four metadata files (questions, lectures,
payments, coupons) provide the content side of the dataset.

\begin{table}[H]
\centering
\small
\begin{tabular}{lll}
\toprule
Column & Type & Description \\
\midrule
\texttt{timestamp}    & int64    & Unix ms; shifted from real values for privacy \\
\texttt{action\_type} & category & 13 types (enter, respond, play\_audio, \dots) \\
\texttt{item\_id}     & string   & Question / bundle / lecture / payment / coupon \\
\texttt{cursor\_time} & float64  & Media playback position in ms \\
\texttt{source}       & category & 8 in-app contexts (sprint, my\_note, adaptive\_offer, \dots) \\
\texttt{user\_answer} & category & a / b / c / d \\
\texttt{platform}     & category & mobile / web \\
\texttt{user\_id}     & int32    & Anonymized student identifier \\
\bottomrule
\end{tabular}
\caption{Schema of the consolidated interaction log.}
\end{table}

The companion \texttt{questions.csv} (13{,}169 rows) records the question's
TOEIC \texttt{part} (1--7), correct answer, and skill tags;
\texttt{lectures.csv} (1{,}021 rows) records the lecture's part, tags, and
video length. Together they enable the integration step that yields a
correctness label for every \texttt{respond} action.

\subsection{Why this dataset}

EdNet-KT4 is appealing for our project for three reasons. First, its
\textbf{raw scale} (131M rows, 6.4 GB) forces us to use efficient
columnar formats and out-of-core operations from the very first step.
Second, its \textbf{action diversity} — beyond just answers — gives the
preprocessing and feature-engineering stages real material to work on
(time gaps, fatigue proxies, lecture engagement, behavioral signals).
Third, the dataset has a natural \textbf{prediction target}: given a
student's history, can we predict whether their next response will be
correct? This is a well-studied knowledge-tracing problem on which we
can compare classical (tree-based) and sequence (LSTM, 1D-CNN) models
on the same evaluation protocol.
